\documentclass[conference]{IEEEtran}

\usepackage{cite}
\usepackage{amsmath,amssymb}
\usepackage{graphicx}
\usepackage{booktabs}
\usepackage{xcolor}
\usepackage{hyperref}
\usepackage{url}

% Red placeholders for results and screenshots that do not exist yet.
\newcommand{\todo}[1]{\textcolor{red}{\textbf{[TODO: #1]}}}
% Includes figs/<file> when it exists; otherwise shows a red placeholder box.
\newcommand{\figbox}[3]{%
  \IfFileExists{figs/#1}{\includegraphics[width=#2]{figs/#1}}%
  {\fbox{\parbox{0.9\columnwidth}{\centering\todo{#3}}}}}

\begin{document}

\title{Four Generative Models for Image Restoration and Face-to-Sketch Synthesis:\\
Autoencoders, Hard and Soft Routing, and a Style-Conditioned cGAN}

\author{\IEEEauthorblockN{Muhammad Harris}
\IEEEauthorblockA{Student ID i232535\\
Generative AI, Assignment 1\\
\todo{affiliation and email}}}

\maketitle

\begin{abstract}
We study four related generative systems. A universal denoising autoencoder restores
images corrupted by salt-and-pepper noise, Gaussian blur, or rectangular occlusion without
being told which corruption is present. A corruption classifier routes each image to one of
three specialist autoencoders (hard routing), and a soft mixture-of-experts variant blends the
same specialists with learned, differentiable weights. A conditional GAN with a U-Net generator
and a PatchGAN discriminator converts face photographs into sketches in one of three styles.
Every model is tuned with Optuna. The planned system serves all four models through one
browser application with a FastAPI backend and a React frontend, with models exported to ONNX. This report
documents the research behind each architectural choice, the alternatives we rejected and why,
the hyperparameter studies, and the evaluation protocol. \todo{insert final headline results
once all training runs finish}
\end{abstract}

\begin{IEEEkeywords}
denoising autoencoder, mixture of experts, hard routing, conditional GAN, PatchGAN, U-Net,
Optuna, ONNX
\end{IEEEkeywords}

%=====================================================================
\section{Introduction}
Photographs are often damaged by noise, blur, or occlusion, and a common approach is to learn a
single restoration network. A different approach is to first diagnose the damage and then
dispatch the image to a specialist. This report compares the two. We also study whether soft,
differentiable routing can combine specialists better than a hard switch. The fourth task is a
different problem: translating a face photograph into a stylized sketch with a conditional
generative adversarial network (cGAN).

The assignment was designed to test research and engineering together. Each design decision in
this report is justified by the literature and by experiments, and each alternative we considered
is listed with the reason it was not chosen.

%=====================================================================
\section{Related Work}
\textbf{Denoising autoencoders.} Vincent et al.\ showed that training an autoencoder to
reconstruct a clean input from a corrupted version yields representations that are robust to the
corruption~\cite{vincent2008dae}. This is the training objective used in Task~1.

\textbf{Upsampling artifacts.} Transposed convolutions can produce checkerboard patterns because of
uneven kernel overlap. Odena et al.\ show that separating resizing from convolution avoids the
problem~\cite{odena2016checkerboard}. We use this upsample-then-convolve pattern in every decoder.

\textbf{Losses for restoration.} Zhao et al.\ compare pixel losses and perceptual losses and find
that combining L1 with a structural term gives better results than L2 alone~\cite{zhao2017loss}.
The structural similarity index (SSIM) was introduced by Wang et al.~\cite{wang2004ssim}.

\textbf{Normalization.} Batch normalization depends on batch statistics, which become noisy at
small batch sizes. Group normalization removes that dependence~\cite{wu2018gn}. Instance
normalization is common in image-to-image networks~\cite{ulyanov2016in}.

\textbf{Mixture of experts.} Jacobs et al.\ introduced adaptive mixtures of local experts, in
which a gating network assigns weights to experts~\cite{jacobs1991moe}. Shazeer et al.\ scaled this
to large networks~\cite{shazeer2017moe}. A known failure mode is routing collapse, where one
expert receives almost all inputs; Task~3 uses a balance regularizer against it.

\textbf{Conditional image translation.} Isola et al.\ introduced pix2pix, a conditional GAN with a
U-Net generator and a PatchGAN discriminator, for paired image translation~\cite{isola2017pix2pix}.
Our Task~4 model follows this design. Perez et al.\ propose feature-wise linear modulation (FiLM) as
a general conditioning mechanism~\cite{perez2018film}, which we considered and rejected (Section~VI).

\textbf{Hyperparameter optimization.} Optuna uses a define-by-run API and pruning of unpromising
trials~\cite{akiba2019optuna}.

\textbf{Datasets.} The Oxford-IIIT Pet dataset~\cite{parkhi2012pets} provides the images for
Tasks~1--3. The FS2K dataset~\cite{fan2021fs2k} provides the photo-sketch pairs for Task~4.

%=====================================================================
\section{Dataset Preparation}
\subsection{Oxford-IIIT Pet (Tasks 1--3)}
The official training and validation list (\texttt{trainval.txt}, 3{,}680 images) is the
development set. We split it 80/20 with seed 42, giving 2{,}944 training and 736 validation images.
The official test list (\texttt{test.txt}, 3{,}669 images) is used only for evaluation. All images are
converted to RGB and resized to $128\times128$. Corruptions are generated at runtime during training
and stored in fixed manifests for validation and testing. Table~\ref{tab:corruptions} gives the
definitions.

\begin{table}[t]
\centering
\caption{Corruption configuration. Training samples severity at random; test uses the fixed levels.}
\label{tab:corruptions}
\begin{tabular}{@{}llll@{}}
\toprule
Condition & Training & Test low & Test medium / high \\
\midrule
Clean & none & none & none \\
Salt-and-pepper & $p\sim U(0.02,0.15)$ & $p=0.03$ & $0.08$ / $0.15$ \\
Gaussian blur & $k\in\{3,5,7\}$, $\sigma\sim U(0.5,2.5)$ & $(3,0.7)$ & $(5,1.5)$ / $(7,2.5)$ \\
Occlusion & 1--3 boxes, 10--35\% area & 1 box, 10\% & 2 boxes, 20\% / 3 boxes, 35\% \\
\bottomrule
\end{tabular}
\end{table}

\subsection{FS2K (Task 4)}
FS2K contains 2{,}104 photo-sketch pairs in three styles~\cite{fan2021fs2k}. We use the official
training file (1{,}058 pairs) and the official test file (1{,}046 pairs). From the training file we
hold out 15\% for validation (159 pairs), stratified by style with seed 42. The split is 899
training, 159 validation, and 1{,}046 test images. Photos and sketches are resized to $128\times128$,
and augmentation (horizontal flip, random crop) is applied to each pair with the same random draw.
Style labels come from the \texttt{style} field of the annotations, because the folder structure does
not match the styles.

%=====================================================================
\section{Task 1: Universal Denoising Autoencoder}
\subsection{Architecture and alternatives}
Table~\ref{tab:t1choices} lists each design choice, the options considered, and the reason for the
choice. The resulting model is fully convolutional. There are no skip connections, so the bottleneck
must carry the information needed for reconstruction.

\begin{table*}[t]
\centering
\caption{Task 1 design decisions.}
\label{tab:t1choices}
\begin{tabular}{@{}p{2.6cm}p{5.2cm}p{5.2cm}p{4.2cm}@{}}
\toprule
Decision & Options considered & Chosen & Reason \\
\midrule
Downsampling & max pooling; strided convolution & strided conv, $k{=}3$, $s{=}2$ & the network learns the downsampling instead of discarding information with a fixed max \\
Upsampling & transposed conv; upsample + conv & upsample + conv & avoids checkerboard artifacts~\cite{odena2016checkerboard} \\
Normalization & BatchNorm; GroupNorm; none & GroupNorm (8 groups) & batch size is searched over 8--64; GroupNorm does not depend on batch statistics~\cite{wu2018gn} \\
Activation & ReLU; LeakyReLU(0.2) & searched by Optuna & both are cheap; the search decides \\
Bottleneck & flattened dense vector; spatial feature map & spatial feature map & preserves spatial layout; no large dense layer \\
Skip connections & none; limited skips & none & the assignment requires a genuine bottleneck \\
Output & linear; sigmoid & sigmoid & images are in $[0,1]$ \\
Loss & L1; L2; L1+SSIM & $\alpha L_1+(1-\alpha)(1-\mathrm{SSIM})$ & L1 gives accuracy, SSIM gives structure~\cite{zhao2017loss,wang2004ssim} \\
\bottomrule
\end{tabular}
\end{table*}

The bottleneck was tested at two depths and two widths: depth $4$ gives an $8\times8$ map and depth $5$
gives $4\times4$. Channels double at each stage from a base width. The compression ratio is the input
size divided by the bottleneck size, $49{,}152/(H\cdot W\cdot C)$. For example, $8\times8\times256$
gives $3\times$ compression and $8\times8\times512$ gives $1.5\times$.

\subsection{Optuna study}
The study ran 30 trials, each for three epochs on the training split, with median pruning. The
search space was: learning rate in $\{10^{-2},3\times10^{-3},10^{-3},3\times10^{-4},10^{-4}\}$; batch size
in $\{8,16,32,64\}$; spatial size $\{4,8\}$; base channels $\{16,32,64\}$; dropout in $[0,0.3]$;
$\alpha\in[0.5,0.95]$; activation in $\{$ReLU, LeakyReLU$\}$. Table~\ref{tab:t1optuna} lists the
trials that matter for the decision.

\begin{table}[t]
\centering
\caption{Task 1 Optuna trials (validation loss, lower is better). Trial 26 was the unconstrained
winner. Trial 17 was the best configuration with $3\times$ or greater compression.}
\label{tab:t1optuna}
\begin{tabular}{@{}rlrlr@{}}
\toprule
Trial & Bottleneck & Compression & Activation & Val.\ loss \\
\midrule
26 & $8\times8\times512$ & $1.5\times$ & ReLU & 0.0830 \\
15 & $8\times8\times512$ & $1.5\times$ & LeakyReLU & 0.0933 \\
17 & $8\times8\times256$ & $3\times$ & LeakyReLU & 0.0923 \\
3 & $8\times8\times128$ & $6\times$ & LeakyReLU & 0.1501 \\
13 & $4\times4\times256$ & $12\times$ & LeakyReLU & 0.1472 \\
7 & $4\times4\times1024$ & $3\times$ & ReLU & 0.1794 \\
\bottomrule
\end{tabular}
\end{table}

Trial 26 had the lowest validation loss, but its bottleneck compresses the input by only $1.5\times$.
Trial 17 is $0.0923$, about 11\% higher, and compresses by $3\times$. It keeps
$3\times$ compression. We chose trial 17. The loss increase is small, and the bottleneck is a
meaningful compression of the input, which the assignment requires. Trial 13 has stronger
compression ($12\times$) but a loss 59\% higher than trial 17. We judged that too large a cost.

The selected hyperparameters are: learning rate $3\times10^{-4}$, batch size 16, spatial size 8,
base channels 32, dropout $0.158$, $\alpha=0.927$, LeakyReLU.

\subsection{Training and evaluation}
The final model is trained for 50 epochs on the training split with the selected settings. Evaluation
uses the official test list and the fixed test grid: each image is tested with the clean input and
with three severities of each corruption, ten entries per image. Results are reported as L1, PSNR,
and SSIM per corruption and severity, and include twelve examples (clean, corrupted, restored,
absolute error map).

Each entry is scored three ways: the model's restored output, the corrupted input itself (no restoration),
and, for salt-and-pepper only, a $3\times3$ median filter. Gains are the model minus the baseline, so a
positive gain means the model improves on that baseline.

\begin{table*}[t]
\centering
\caption{Task 1 test results on the official test list, 3{,}669 images per row. Gain = model minus
input (or minus median filter, where given). Clean rows have no corruption, so their input equals the
target and gains are not defined.}
\label{tab:t1results}
\small
\begin{tabular}{@{}llrrrrrrrr@{}}
\toprule
 & & \multicolumn{4}{c}{PSNR (dB)} & \multicolumn{4}{c}{SSIM} \\
\cmidrule(lr){3-6}\cmidrule(lr){7-10}
Corruption & Severity & Model & Input & Gain & Median & Model & Input & Gain & Median \\
\midrule
Clean & none & 23.82 & -- & -- & -- & 0.696 & 1.000 & -- & -- \\
Salt-and-pepper & low & 23.72 & 20.15 & +3.57 & 32.19 & 0.689 & 0.613 & +0.076 & 0.926 \\
Salt-and-pepper & medium & 23.47 & 15.89 & +7.59 & 31.12 & 0.676 & 0.356 & +0.319 & 0.918 \\
Salt-and-pepper & high & 23.03 & 13.16 & +9.87 & 28.99 & 0.657 & 0.224 & +0.434 & 0.899 \\
Blur & low & 23.89 & 34.20 & $-$10.31 & -- & 0.695 & 0.962 & $-$0.266 & -- \\
Blur & medium & 23.90 & 28.14 & $-$4.24 & -- & 0.690 & 0.855 & $-$0.165 & -- \\
Blur & high & 23.68 & 25.37 & $-$1.69 & -- & 0.676 & 0.755 & $-$0.080 & -- \\
Occlusion & low & 22.10 & 16.78 & +5.32 & -- & 0.665 & 0.880 & $-$0.215 & -- \\
Occlusion & medium & 20.71 & 13.67 & +7.04 & -- & 0.636 & 0.777 & $-$0.142 & -- \\
Occlusion & high & 19.06 & 11.48 & +7.58 & -- & 0.596 & 0.653 & $-$0.057 & -- \\
\bottomrule
\end{tabular}
\end{table*}

The results show three distinct patterns. For salt-and-pepper noise, the model improves on the noisy input
by up to 9.9\,dB, but a simple median filter is better by 5 to 8\,dB at every severity. For blur, the
model's output scores below the blurred input at every severity, so restoration makes the image worse. For
occlusion, the model improves on the input by 5 to 8\,dB, but its output reaches only 19 to 22\,dB. Even on clean
images, the model reaches only 23.8\,dB and an SSIM of 0.70, which means its output is smooth even when
nothing needs fixing.

\figbox{t1_training_curve.png}{\columnwidth}{Task 1 training and validation loss curves}

\figbox{t1_examples.png}{\columnwidth}{Task 1 examples: clean, corrupted, restored, error map}

\subsection{Failure cases}
We identify four failure cases from the test results. Each is backed by the numbers above. The
explanations are hypotheses, and the example images below should be inspected to
confirm them.

\textbf{Failure 1: the model ignores its blurred input.} At low blur severity, the model scores 23.9\,dB,
while the blurred input scores 34.2\,dB. The model's output looks much the same regardless of input,
which suggests it has learned a generic smooth image. Our hypothesis is that the bottleneck and the
training objective favour smooth outputs, so the model discards the fine detail that a mild blur leaves
intact.

\textbf{Failure 2: noise removal is weaker than a median filter.} At high salt-and-pepper severity, the
model reaches 23.0\,dB, a gain of 9.9\,dB over the noisy input, but the median filter reaches 29.0\,dB. A
local median filter handles impulse noise directly, which the learned model does not match. We
hypothesise that the L1 and SSIM loss mix does not penalise the model enough for keeping isolated noise
pixels.

\textbf{Failure 3: the reconstruction ceiling.} On clean images, the model reaches only 23.8\,dB and an
SSIM of 0.70. Every corrupted entry is restored to at most this quality, so the ceiling limits all the
other results. We hypothesise that the 8$\times$8$\times$256 bottleneck (3$\times$ compression) discards
fur texture and edges that a larger bottleneck would keep. The Optuna search favoured the 1.5$\times$
bottleneck on loss alone (Table~\ref{tab:t1optuna}), which points the same way.

\textbf{Failure 4: large occlusions cannot be filled.} At high occlusion severity, the input covers up to 35\%
of the image, and the model reaches 19.1\,dB and SSIM 0.60. The model can't reconstruct content hidden
behind a large black box from the compressed code. This is the one corruption where restoration improves
on the input, but the result is still far below the clean reconstruction.

%=====================================================================
\section{Task 2: Corruption Classification and Hard-Routed Specialists}
\subsection{Classifier}
The classifier uses the same down-block as Task~1 (strided convolution, GroupNorm, activation), four
stages, then global average pooling and a linear layer with four outputs. We chose global average
pooling over flattening because the decision is whether a corruption is present anywhere, not where
it is. Batches are balanced across the four classes with a custom sampler, so the classifier cannot
learn a bias toward one class.

The Optuna search tuned learning rate, batch size, base channels, dropout, and weight decay. The
selected values are: learning rate $3\times10^{-4}$, batch size 16, base channels 32, dropout
$0.018$, weight decay $2.0\times10^{-5}$. \todo{add the classifier trial table; the run log was not
saved to a file}

\subsection{Specialists}
All three specialists use the same autoencoder as Task~1. The architecture was chosen with one shared
search on pooled corrupted data, then three copies were trained independently, one per corruption.

The shared search ran 20 trials of three epochs each. The unconstrained best trial (\#10) had a
$8\times8\times512$ bottleneck, which compresses only $1.5\times$. Within the four bottleneck options we
set in the design, the best trial was \#12, with an $8\times8\times256$ bottleneck and validation loss
$0.1196$. We selected trial 12. It is 32\% higher in loss than trial 10, and it keeps the
$3\times$ compression required by the design. Table~\ref{tab:t2spec} lists the selected configuration.

\begin{table}[t]
\centering
\caption{Task 2 specialist search, selected configuration (trial 12).}
\label{tab:t2spec}
\begin{tabular}{@{}lr@{}}
\toprule
Parameter & Value \\
\midrule
Bottleneck & $8\times8\times256$ \\
Base channels & 32 \\
Learning rate & $3\times10^{-4}$ \\
Batch size & 32 \\
Dropout & $0.280$ \\
$\alpha$ & $0.870$ \\
Activation & LeakyReLU \\
Validation loss (pooled) & $0.1196$ \\
\bottomrule
\end{tabular}
\end{table}

\subsection{Routing and evaluation}
At inference the classifier chooses an expert. Clean predictions bypass the experts. We evaluate two
modes. In \emph{oracle routing} the true label chooses the expert, which isolates specialist quality.
In \emph{predicted routing} the classifier chooses, which measures the full system. The difference
between the two is the cost of classification errors. We also collect the images where a
misclassification most increased the error relative to oracle routing.

\todo{classifier test metrics (accuracy, macro P/R/F1, per-class, normalized confusion matrix)}
\todo{oracle versus predicted routing table per corruption and severity}
\figbox{t2_confusion.png}{0.8\columnwidth}{normalized four-class confusion matrix}
\figbox{t2_failures.png}{\columnwidth}{classifier-caused failure cases}
\figbox{t2_classifier_curve.png}{0.8\columnwidth}{classifier training loss and validation accuracy}
\figbox{t2_specialist_curves.png}{0.8\columnwidth}{specialist validation loss per corruption}

%=====================================================================
\section{Task 3: Soft Mixture-of-Experts Restoration}
\subsection{Architecture and alternatives}
Hard routing sends each image to one expert and cannot handle mixed corruptions or uncertain
classifications. We replace it with a soft mixture: a gating network produces weights
$w=\mathrm{softmax}(G(\tilde x)/\tau)$ over the identity branch and the three experts, and the output is
$\hat x=\sum_k w_k A_k(\tilde x)$. Table~\ref{tab:t3choices} lists the alternatives.

\begin{table}[t]
\centering
\caption{Task 3 design decisions.}
\label{tab:t3choices}
\begin{tabular}{@{}p{2.2cm}p{2.9cm}p{2.6cm}@{}}
\toprule
Decision & Options considered & Chosen \\
\midrule
Gate & new network; Task~2 classifier & Task~2 classifier \\
Initialization & random; pretrained & pretrained from Task~2 \\
Training & joint from start; warm-up then joint & warm-up (experts frozen), then joint \\
Routing & hard; softmax with temperature & softmax with temperature $\tau$ \\
Collapse guard & none; balance loss & balance loss $\sum_k(\bar w_k-\tfrac14)^2$ \\
Loss weight $\lambda_s$ & tuned freely; $1-\lambda_1$ & $1-\lambda_1$ (keeps weights complementary) \\
\bottomrule
\end{tabular}
\end{table}

\subsection{Training and Optuna}
Training has two stages. In warm-up the experts are frozen and only the gate learns, so it learns
routing against good experts. In joint fine-tuning everything is unfrozen at a lower learning rate. The
loss is $\mathcal{L}=\lambda_1L_1+\lambda_s(1-\mathrm{SSIM})+\lambda_cL_{CE}+\lambda_bL_{\mathrm{balance}}$,
starting from $0.8/0.2/0.1/0.01$. The Optuna search covers joint learning rate in $[10^{-6},10^{-4}]$, temperature
in $[0.3,3]$, $\lambda_c\in[0.01,0.5]$, $\lambda_b\in[0.001,0.1]$, and $\lambda_1\in[0.5,0.9]$. Trials are pruned
when two or more experts become inactive. \todo{insert Optuna study results: trial count, best trial,
selected configuration, and pruning count}

\subsection{Routing analysis}
We report the mean routing weights for each true corruption type and severity, and check for dead
experts (mean weight below 0.05) and for experts that dominate unrelated inputs. \todo{routing table
per corruption and severity; routing heatmap; examples where one expert dominates and where weights
are spread}
\figbox{t3_routing_heatmap.png}{0.9\columnwidth}{routing weights by corruption and severity}
\figbox{t3_routing_examples.png}{\columnwidth}{routing examples: one expert dominates (top rows) versus weights spread (bottom rows); columns are clean, corrupted, restored}
\figbox{t3_training_curve.png}{0.8\columnwidth}{soft mixture joint fine-tuning validation loss}

\subsection{Results}
\todo{reconstruction results of the soft mixture compared with Task 1 and Task 2 on the official test set}

%=====================================================================
\section{Task 4: Style-Conditioned Face-to-Sketch cGAN}
\subsection{Style conditioning: alternatives and choice}
Three ways of conditioning the generator on style were considered. \emph{(A) Input concatenation:} a
learned embedding for each style is repeated over the image and concatenated to the input, in both the
generator and the discriminator. \emph{(B) Bottleneck injection:} the embedding is applied only at the
U-Net bottleneck. \emph{(C) FiLM:} feature-wise scaling and shifting at each stage~\cite{perez2018film}.

We chose (A). Its main weakness is that style information reaches the decoder only through skip
connections, so the decoder may under-use it. B addresses that but adds a second conditioning
point. C is the most expressive, but with three styles and 899 training images the added machinery is
hard to justify. \todo{compare A and B on validation; if B is clearly better, revise this decision}

\subsection{Architecture and alternatives}
The generator is a U-Net with five downsampling stages, from $128$ to a $4\times4$ bottleneck, so the
skip connections keep detail at $64$, $32$, $16$, and $8$ pixels. Decoders use upsample-then-convolve
and InstanceNorm, with dropout in the first three blocks. The output is tanh on $[-1,1]$. We chose
InstanceNorm over BatchNorm because batch sizes from 8 to 32 are searched. The discriminator is a
PatchGAN with two stride-2 layers and two stride-1 layers. Its receptive field is 34 pixels, which
matches the relative patch size of pix2pix at this resolution~\cite{isola2017pix2pix}. Table~\ref{tab:t4choices}
summarizes the remaining options.

\begin{table}[t]
\centering
\caption{Task 4 design decisions.}
\label{tab:t4choices}
\begin{tabular}{@{}p{2.2cm}p{2.9cm}p{2.6cm}@{}}
\toprule
Decision & Options considered & Chosen \\
\midrule
Generator & encoder-decoder; U-Net & U-Net (skips needed for edges) \\
Depth & 4, 5, 6 stages & 5 stages (4$\times$4 bottleneck) \\
Normalization & BatchNorm; InstanceNorm & InstanceNorm \\
Discriminator & whole-image; PatchGAN & PatchGAN, RF 34 px \\
Adversarial loss & BCE; least squares & BCE with logits (assignment) \\
Output & sigmoid; tanh & tanh on $[-1,1]$ \\
\bottomrule
\end{tabular}
\end{table}

\subsection{Loss and training}
The generator minimizes $\mathcal{L}_G=\mathcal{L}_{adv}+\lambda_{L1}\mathcal{L}_{L1}$ and the
discriminator minimizes binary cross-entropy on real and generated pairs. We log the discriminator real
loss, discriminator fake loss, generator adversarial loss, generator L1 loss, and validation L1 and SSIM.
Optuna searches generator and discriminator learning rates ($[10^{-4},4\times10^{-4}]$, log scale), batch
size $\{8,16,32\}$, base channels $\{32,64\}$, dropout $[0,0.5]$, style embedding size $\{8,16,32\}$, and
$\lambda_{L1}\in[10,200]$. Trials run 15 epochs; the winner is retrained for 200 epochs.
The search ran 20 trials of 15 epochs each, using the full official training split (899 training and 159
validation images). The test set was not used. Median pruning stopped 6 of the 20 trials. Trials 0--5 of
the study failed on a filename error before any training, so they are excluded. Table~\ref{tab:t4optuna}
lists the five best trials.

\begin{table}[t]
\centering
\caption{Task 4 Optuna results: five best trials by validation L1 (lower is better).}
\label{tab:t4optuna}
\small
\begin{tabular}{@{}rrrrrrrrr@{}}
\toprule
Trial & Val L1 & $\mathrm{lr}_G$ & $\mathrm{lr}_D$ & Batch & Base & Dropout & Style dim & $\lambda_{L1}$ \\
\midrule
25 & 0.0928 & $2.3\times10^{-4}$ & $3.5\times10^{-4}$ & 8 & 64 & 0.29 & 32 & 199.7 \\
23 & 0.0934 & $1.5\times10^{-4}$ & $3.1\times10^{-4}$ & 8 & 64 & 0.48 & 32 & 193.6 \\
21 & 0.0935 & $1.8\times10^{-4}$ & $2.9\times10^{-4}$ & 8 & 64 & 0.38 & 32 & 173.8 \\
16 & 0.0947 & $3.4\times10^{-4}$ & $2.4\times10^{-4}$ & 8 & 64 & 0.37 & 32 & 132.9 \\
24 & 0.0952 & $1.1\times10^{-4}$ & $3.4\times10^{-4}$ & 8 & 64 & 0.44 & 32 & 163.2 \\
\bottomrule
\end{tabular}
\end{table}

The five best trials are within 0.0024 of each other, which is within the noise of a single seed, so the
ranking among them is not reliable. All five use batch size 8, base width 64, and style embedding 32. The
best $\lambda_{L1}$ (199.7) sits at the top of its range, so the true optimum may lie higher. The selected
configuration is trial 25, and it is used for both conditioning options.

\subsection{Results}
\todo{validation and test L1 and SSIM per style}
\figbox{t4_samples.png}{\columnwidth}{photo, generated sketch, and true sketch for fixed validation photos at several epochs}
\figbox{t4_loss_curves.png}{\columnwidth}{discriminator and generator losses (left) and L1 (right) during training}
\figbox{t4_conditioning_comparison.png}{0.8\columnwidth}{validation L1 for option A (input) versus option B (bottleneck)}

%=====================================================================
\section{Application Architecture}
The four models are served by one FastAPI backend, which validates uploads, preprocesses images, runs
the ONNX models, and returns outputs with routing information and timing. The React frontend provides
four workspaces: Universal Restoration, Hard-Routed Restoration, Soft Mixture-of-Experts Restoration, and
Face-to-Sketch Generator. The interface was designed in Google Stitch before it was built, and
Figure~\ref{fig:stitch} shows the four screens. The implemented endpoints are \texttt{/health},
\texttt{/samples}, \texttt{/universal-restoration}, \texttt{/hard-routing}, \texttt{/soft-mixture}, and
\texttt{/face-to-sketch}. Each returns the restored or generated image as PNG with its inference time. An
endpoint whose model is not loaded returns HTTP 503 rather than failing.

\begin{figure*}[t]
\centering
\figbox{stitch_universal.png}{0.48\textwidth}{Stitch design: Universal Restoration}
\hfill
\figbox{stitch_hard_routed.png}{0.48\textwidth}{Stitch design: Hard-Routed Restoration}

\vspace{0.5em}
\figbox{stitch_soft_moe.png}{0.48\textwidth}{Stitch design: Soft Mixture-of-Experts}
\hfill
\figbox{stitch_face_sketch.png}{0.48\textwidth}{Stitch design: Face-to-Sketch}
\caption{Google Stitch designs for the four workspaces. The implemented interface follows these layouts.}
\label{fig:stitch}
\end{figure*}

\todo{application screenshots, after the app runs; inference times measured on the target machine}
\todo{Docker Compose: describe the compose file and the one-command start, once it is written}

%=====================================================================
\section{Limitations}
\begin{itemize}
\item Optuna searches used short trials of two to three epochs, so the chosen hyperparameters are
      estimates. Full-schedule results may differ.
\item Each experiment was run once with one seed. Differences of a few percent should not be
      over-interpreted.
\item The GAN training may show decreasing losses while sample quality worsens. We rely on
      validation L1 and SSIM plus visual inspection.
\item FS2K is small (899 training pairs), which limits how much a GAN can learn.
\end{itemize}

%=====================================================================
\section{Conclusion}
\todo{summarize the findings after all evaluations are complete}

%=====================================================================
\section*{Reproducibility}
Source code: \todo{GitHub repository URL}. Demonstration video (YouTube, unlisted):
\todo{YouTube URL}. Model weights are available through the download links in the repository README.

%=====================================================================
\appendix
\section*{AI-Use Appendix}
\begin{tabular}{@{}p{3.2cm}p{4.6cm}p{4.6cm}@{}}
\toprule
Tool & Used for & How outputs were checked \\
\midrule
Claude (Anthropic) & research discussion, code drafts, debugging, report drafting & every script was run on the
real dataset; numbers were taken from logs; citations need checking before submission \\
\bottomrule
\end{tabular}

%=====================================================================
\begin{thebibliography}{99}
\bibitem{vincent2008dae} P.~Vincent, H.~Larochelle, I.~Lajoie, Y.~Bengio, and P.-A.~Manzagol,
``Extracting and composing robust features with denoising autoencoders,'' in \emph{Proc. ICML}, 2008.
\bibitem{odena2016checkerboard} A.~Odena, V.~Dumoulin, and C.~Olah, ``Deconvolution and checkerboard
artifacts,'' \emph{Distill}, 2016.
\bibitem{zhao2017loss} H.~Zhao, O.~Gallo, I.~Frosio, and J.~Kautz, ``Loss functions for image restoration
with neural networks,'' \emph{IEEE Trans. Computational Imaging}, 2017.
\bibitem{wang2004ssim} Z.~Wang, A.~C.~Bovik, H.~R.~Sheikh, and E.~P.~Simoncelli, ``Image quality assessment:
from error visibility to structural similarity,'' \emph{IEEE Trans. Image Processing}, 2004.
\bibitem{wu2018gn} Y.~Wu and K.~He, ``Group normalization,'' in \emph{Proc. ECCV}, 2018.
\bibitem{ulyanov2016in} D.~Ulyanov, A.~Vedaldi, and V.~Lempitsky, ``Instance normalization: The missing
ingredient for fast stylization,'' \emph{arXiv:1607.08022}, 2016.
\bibitem{jacobs1991moe} R.~A.~Jacobs, M.~I.~Jordan, S.~J.~Nowlan, and G.~E.~Hinton, ``Adaptive mixtures of
local experts,'' \emph{Neural Computation}, 1991.
\bibitem{shazeer2017moe} N.~Shazeer et al., ``Outrageously large neural networks: The sparsely-gated
mixture-of-experts layer,'' in \emph{Proc. ICLR}, 2017.
\bibitem{isola2017pix2pix} P.~Isola, J.-Y.~Zhu, T.~Zhou, and A.~A.~Efros, ``Image-to-image translation with
conditional adversarial networks,'' in \emph{Proc. CVPR}, 2017.
\bibitem{perez2018film} E.~Perez, F.~Strub, H.~de~Vries, V.~Dumoulin, and A.~Courville, ``FiLM: Visual
reasoning with a general conditioning layer,'' in \emph{Proc. AAAI}, 2018.
\bibitem{akiba2019optuna} T.~Akiba, S.~Sano, T.~Yanase, T.~Ohta, and M.~Koyama, ``Optuna: A next-generation
hyperparameter optimization framework,'' in \emph{Proc. KDD}, 2019.
\bibitem{parkhi2012pets} O.~M.~Parkhi, A.~Vedaldi, A.~Zisserman, and C.~V.~Jawahar, ``Cats and dogs,'' in
\emph{Proc. CVPR}, 2012.
\bibitem{fan2021fs2k} D.-P.~Fan et al., ``FS2K: Facial sketch synthesis dataset,'' 2021. \todo{verify citation and link}
\end{thebibliography}

\end{document}
